% !TEX program = pdflatex
%
% Companion to judgements-and-warrants.tex: what the branch
% `numeric-core-and-verification-judgement` actually built, where it departs from the proposal,
% and what the proposal specifies that is not built.
%
% Build:  pdflatex judgements-and-warrants-implementation.tex  (twice, to resolve references)

\documentclass[manuscript,nonacm=true]{acmart}

\setcopyright{none}
\settopmatter{printacmref=false, printccs=false, printfolios=true}

\usepackage{amsmath}
\usepackage{booktabs}

\setlength{\parindent}{0pt}
\setlength{\parskip}{0.30\baselineskip plus 0.15\baselineskip minus 0.1\baselineskip}

\newcommand{\code}[1]{\texttt{\small #1}}
\newcommand{\sem}[1]{[\![#1]\!]}
\newcommand{\Prop}{\code{Prop}}
\newcommand{\cl}[1]{\textbf{#1}}

\begin{document}

\title[Judgements and Warrants: Implementation]{Judgements, Warrants, and Logics\\[0.5em]
\large Implementation companion}
\author{Hans-Martin Will}
\affiliation{%
  \institution{The Eigenius Project}
  \country{USA}
}
\email{hwill@acm.org}
\renewcommand{\shortauthors}{Will}
\date{September 2026}

\begin{abstract}
The proposal \emph{Judgements, Warrants, and Logics} opens by stating that its design is unimplemented. That is no longer accurate. This companion records what was built, names each construct's location in the source, and separates three things the proposal's own text cannot distinguish: what is implemented and load-bearing, what is implemented differently from the proposal and why, and what the proposal specifies that does not exist. It also records the additions the proposal does not describe, which arose from making the \emph{Verified} ground reachable for a numeric claim.
\end{abstract}

\maketitle

\section*{Status}

The proposal's Status section reads \emph{``this is a proposal, and the design is currently unimplemented.''} The design is implemented. The two strata, the three grounds, the witness family, the warrant algebra, the well-foundedness condition and the uniform check-mode rule are all in the kernel and exercised by end-to-end tests. That sentence in the proposal should be revised rather than left to mislead a reader who arrives at it later.

This document is written against the branch \code{numeric-core-and-verification-judgement}. Source references are paths in the Eigenius repository and are accurate at the time of writing; they will drift, and the construct names are the durable part.

\section{Correspondence}

Each row names a construct in the proposal and where it lives.

\begin{table}[h]
\small
\begin{tabular}{@{}p{0.30\linewidth}p{0.62\linewidth}@{}}
\toprule
\textbf{proposal} & \textbf{implementation} \\
\midrule
$\mathrm{Judgement}(L, t, T)$ &
\code{eigentt:Judgement}, constructor \code{holds(logic, term, type)}, declared in the core ontology \\

$\mathrm{JustifiedBy}(j, P)$ &
\code{justification:Certificate : Prop -> Type 2}, seven constructors, in \code{ontologies/justification/justification.esl} \\

three grounds &
\code{declared} / \code{observed} / \code{verified} constructors, each consuming a witness \\

witness family &
\code{witness:IsDeclaredAs}, \code{IsObservedAs}, \code{IsVerifiedAs} --- zero-constructor inductives in \Prop{}, declared in \emph{core} because the kernel inhabits them \\

\emph{Computed} / \emph{Sampled} &
not stored, and no enumeration names them; they are read off term shape by the projection algebra \\

warrant computed, never stored &
\code{support}, \code{is\_fully\_verified}, \code{leaves\_of}, \code{survives\_without}, \code{cited\_iris} in \code{kernel/src/justification/} \\

well-foundedness &
\code{kernel/src/justification/wellfounded.rs}, over the \code{core:mentions} index \\

provenance as relations &
\code{prov:was\_attributed\_to}, \code{was\_generated\_by}, \code{used}, \code{had\_primary\_source}, all resource-typed \\

verification institution &
\code{crates/eigenius-lean}, wrapping \code{nanoda\_lib} in-process \\

statistics institution &
\code{crates/eigenius-statistics} \\
\bottomrule
\end{tabular}
\end{table}

\section{Where the implementation departs from the proposal}

\cl{The justification family is indexed by the proposition alone.} The proposal writes $\mathrm{JustifiedBy}(j, P)$, naming a relation between a term and a proposition. The implementation declares \code{Certificate : Prop -> Type 2}: the certificate \emph{is} the term, so an inhabitant of \code{Certificate(P)} is a constructor tree and the family carries one index. Nothing is lost. The projection algebra walks the value where it would otherwise walk the index; two derivations of one $P$ remain distinct values; and the proposal's own primitive is the typing $t : F$, which is this shape. \code{JustifiedBy(j, P)} in the stratification table names the relation, not an encoding of it.

\cl{Seven constructors, not six.} The proposal's Appendix counts the algebra down to six. The implementation has seven, and the two extra facts are these. \emph{Sum is split} into \code{sum\_l} and \code{sum\_r}, which share a type; the constructor chosen is what records which branch the author preferred, and that record has nowhere else to live once the term is the value. \emph{Specialisation is a constructor}: \code{spec\_poly} narrows a universal proposition to an instance, adding no ground, which is what lets one declared bridge serve many claims --- necessary because the algebra has no implication introduction, so a bridge cannot be derived and must be declared once and instantiated.

\cl{Sum requires both branches to be justified.} This departs from LP's axiom, under which a sum needs only one. The reason is the projection: \code{support} reads Sum disjunctively, so an unjustified branch would answer the counterfactual \emph{would this survive losing $X$?} with a ground nobody established --- and it would answer in the reassuring direction.

\cl{The certificate inhabits \code{Type 2}.} A constructor argument's sort may not exceed the inductive's own. \code{spec\_poly} binds \code{T : Type 1}, which places that argument at \code{Sort 3}, so the family sits at \code{Type 2}. This is a consequence of admitting specialisation and has no effect on values.

\section{What is implemented and load-bearing}

\cl{The two strata do not connect, and the enforcement is a refusal.} \code{eigentt:Term} carries a \code{Checked(payload\_iri)} constructor naming a proof an external checker verified. The kernel's \code{check} \emph{rejects} it (\code{kernel/src/nbe/check/mod.rs}): it holds no proof of the offered proposition and will not manufacture one. A hand-authored \code{holds(logic\_lean4, Checked(a), P)} therefore fails at commit, while an institution-emitted one never reaches validation, because \code{structural\_validate} runs before \code{autoonload\_dispatch} and the followup pipeline slice has no validation phase. The asymmetry falls out of the pipeline's shape rather than from a guard. Without it, \emph{Verified} would be assertable by anyone who could write a judgement.

\cl{\emph{Verified} is computed; \emph{Declared} and \emph{Observed} are postulated.} Witness admission (\code{kernel/src/layer/witness\_admission.rs}) maps a resource's provenance trace to a witness category: \code{prov:DeclarationTrace} to \emph{Declared}, \code{prov:ObservationTrace} to \emph{Observed}, \code{prov:VerificationTrace} to \emph{Verified}. The first two are the constant specification the proposal names as a trusted-computing-base element. The third is different in kind: it keys off the trace's own \code{prov:judgement} --- the checker's retained result, \code{holds(logic, t, P)} --- rather than off the target resource's stored proposition, so the ground rests on what was checked and not on what the resource claims about itself.

\cl{A program trace grounds nothing.} \code{prov:ProgramTrace} maps to no witness category. A computed claim does not rest on the fact that a computation ran; it rests on the declaration that the plan denotes a function $I \to O$, which no execution establishes, and on the observed inputs. The composite is built from those two, so the run record is provenance and not a third ground. An author's transcription of a run that happened elsewhere maps to \emph{Declared}, since a transcription establishes only that someone asserts the run occurred.

\cl{The warrant is a function of the retained term.} No resource carries a grade, and the contradictory state --- a resource labelled \emph{Verified} with no proof --- is unrepresentable rather than excluded by a validation rule.

\cl{Well-foundedness is enforced at commit} over the \code{core:mentions} index, with the declared-premise carve-out the proposal argues for: a declared premise has no support graph to inspect, so the condition holds vacuously on it.

\cl{The trusted computing base is inventoried on the chain.} A verification trace records the export bytes, the checked declaration name, the proposition, the checker identity, and --- \code{prov:permitted\_axioms} --- the axiom allowlist as the check actually ran. The allowlist is part of every verdict's trust basis and was previously invisible, so a proof leaning on \code{Classical.choice} and a constructive one produced indistinguishable traces.

\section{Specified and not built}

\cl{\code{justification:proof} is declared and populated nowhere.} It is the optional second judgement \code{holds(logic, t, P)} --- a proof of the conclusion's own proposition, which is the proposal's \emph{Verified} configuration with no middle stratum. The declaration is correct and nothing writes one yet.

\cl{The reverse comorphism does not exist.} The Lean-to-chain direction is live and is what makes \emph{Verified} reachable. The chain-to-Lean direction needs an EigenTT $\to$ Lean term/type translation that D30 does not specify, and the validator rejects a stub authoring, so its absence is the correct state rather than an oversight.

\cl{Warrant is not expressible in the query language.} A justification term lands as an opaque JSON value; the query language's equality has no arm for it, so no pattern binds one and no built-in decodes one. The projections are a Rust API. \emph{Warrant becomes a query} must not be read as \emph{an EigenQL query} --- the \(\lambda\)Prolog extension that would make it one is tracked separately. Provenance, by contrast, \emph{is} a query, which is the point of having made it relational.

\cl{The numeric ceiling is real and not yet lifted.} See below.

\section{Additions the proposal does not describe}

The proposal establishes when a claim \emph{can} reach \emph{Verified}. Making a numeric claim actually reach it required work the proposal does not anticipate.

\cl{A numeric primitive core.} Chain relations over measured quantities are mapped to Lean propositions by a fixed table in the externalizer: $\le$ and IEEE equality are asserted correspondences, and $<$, $>$, $\ge$ derive from them. Each asserted correspondence is a trusted-computing-base entry, so the table is in reviewed Rust and cannot be extended from the chain --- stating it as a chain property would make a TCB entry authorable by committing a resource.

\cl{Float literals are normalised to the form the elaborator produces.} On the pinned toolchain Lean's \code{Float} is opaque: no kernel reduces it, so two encodings of one double are definitionally distinct unless syntactically identical. Both sides must therefore agree on a form, and the externalizer renders a literal as a person writes it. This is version-scoped, not permanent.

\cl{The Lean version is gated.} An export declares the Lean that produced it, and the implementation refuses one outside the supported series, because what \code{Float} \emph{means} changed in Lean 4.33: \code{Float} became a structure over a bit-pattern model, and the ordering relations changed signature from \Prop{} to \code{Bool}. The terms the externalizer builds still elaborate under both, so the skew would otherwise surface only as an unexplained definitional-equality failure. This gate applies where a proposition is externalized and not to a name-level check, since only externalization fixes what Lean constants mean.

\cl{An exact type is the pivot for procedural proofs.} The argument that originally chose \code{Float} was that a proof about a bespoke relation inherits none of Lean's order lemmas. Because the kernel cannot evaluate floats, there are almost no \code{Float} order lemmas to inherit, so that argument fails in both directions. Threshold reasoning --- weakening, and multiple-testing corrections --- needs order lemmas on the type, not evaluation of the term, which makes an exact rational affordable: the statistics institution already abstracts the dataset away, so only thresholds and bounds appear in propositions. If that pivot is taken, a stored double must map to its \emph{exact} binary64 value and not to the decimal its author typed, or a proof can succeed in Lean where the host computation failed.

\section{Evidence}

Two end-to-end tests exercise the design rather than its parts.

\cl{Three grounds compose in one certificate.} A single certificate of the form \code{app(app(declared(rule), verified(claim)), observed(intake))} commits and projects, with a near-miss variant refused. This is the first place \code{verified} appears as an argument to \code{app} --- before it, the composition operators had never been applied to a checked ground.

\cl{Proving one step removes one declaration.} One proposition is authored twice, differing in a single leaf: asserted in one, discharged by a Lean proof in the other. Both commit; the proved variant's declared-ground set is smaller by exactly that leaf and its verified-ground set contains it. Neither is fully verified, and the test asserts that explicitly --- the domain bridge is declared either way, and a conclusion resting on a declared bridge is not fully verified however much below it is proved. What moves is what nobody proved.

\end{document}
